% Copyright (c) 2026 Eugene Joseph M. Ragasa.
% SPDX-License-Identifier: Apache-2.0
% Author: Eugene Joseph M. Ragasa
% Formal proof completed: 2026-10-26
\documentclass[11pt]{article}
\usepackage[T1]{fontenc}
\usepackage{amsmath,amssymb}
\usepackage{booktabs}
\usepackage[margin=0.85in]{geometry}
\usepackage[hidelinks]{hyperref}
\usepackage{tabularx}
\usepackage{titlesec}
\newcolumntype{Y}{>{\raggedright\arraybackslash}X}
\titlespacing*{\section}{0pt}{1.6ex plus .4ex minus .2ex}{0.7ex}

\title{Proof Dossier: Resolvent-Constructed Rank-One Defect Equation}
\author{Eugene Joseph M. Ragasa}
\date{Formal proof completed: 2026-10-26}

\begin{document}
\raggedright
\maketitle

\section{Identity and status}
\begin{tabularx}{\textwidth}{@{}lY@{}}
\toprule
Statement ID & \texttt{QM-RESOLVENT-THM-013} \\
Status & proved; no admissions \\
Lean declaration & \path{ProjectKoios.QuantumMechanics.FiniteDimensional.rankOneDefectMatrix_mulVec_rankOneResolventVector_eq_smul} \\
Owning module & \path{ProjectKoios/QuantumMechanics/FiniteDimensional/RankOneResolvent.lean} \\
Author and license & Eugene Joseph M. Ragasa; \texttt{Apache-2.0} \\
\bottomrule
\end{tabularx}

\section{Exact mathematical statement}
Let $H$ be finite, and suppose $A=H-EI$ is invertible and
$1-g\langle v|A^{-1}|v\rangle=0$.  For $\psi=A^{-1}v$,
\[
 (H-g|v\rangle\langle v|)\psi=E\psi.
\]

\section{Immutable Formal provenance}
The declaration was introduced at Formal commit
\texttt{c516cc46d5d118ca1d316a364e71692cbe892d5e} and merged through PR~\#2 at
\texttt{265fb7cb1a0b448cf3d069d4c1de9e5af1a912f8}.  It is present at the dossier base
\texttt{d1cc9d5f6ed5b5b1cc8e34ad69b3d40d3104c3f4}.  The owning module SHA-256 is
\texttt{4469d83d3b267574b2ec533273a81267b8ed889eeaf70decd1997c0d8e6582cb}; the exact locator is \texttt{QM-RESOLVENT-THM-013}.

The declaration is source-neutral finite rank-one matrix algebra.  A KSD calculation motivated the module, but no authored matrix, vector, energy, root, tolerance, eigensolver output, or retained artifact is bound by this theorem.

\section{Assumptions, conventions, and statement delta}
The shifted matrix invertibility and exact secular-zero premise are supplied.  Energy and magnitude may be complex.  The theorem does not assume $v\ne0$ and therefore does not itself assert $\psi\ne0$.
The Lean proposition is authoritative.  This prose restatement neither strengthens the hypotheses
nor supplies an application binding.

\section{Statement-specific prior art and contribution boundary}
The resolvent construction for a rank-one secular equation is standard \texttt{background-only} finite-dimensional algebra.  The declaration records the exact conditional eigenvalue equation without promoting it to an existence theorem.

\section{Proof structure and mathematical argument}
Apply the kernel-vector theorem to $A=H-EI$: the secular premise yields
$(A-g|v\rangle\langle v|)\psi=0$.  Expand $A$, move $-E\psi$ to the other side, and recognize the remaining matrix as $H-g|v\rangle\langle v|$.

\section{Lean realization and dependencies}
The proof depends on \texttt{rankOneDefectMatrix\_mulVec\_rankOneResolventVector\_eq\_zero}, finite matrix-vector algebra, and exact expansion of the shifted matrix.  The declaration was rechecked with Lean \texttt{4.35.0} at commit \texttt{eb6cafb7290e8bd5f2937522d8d804f58f02bafc} and Mathlib \texttt{56a005855ea17e0b0958f73ec801ae140b0a1679}.

\section{Trusted base and axiom report}
\begin{verbatim}
#print axioms
  ProjectKoios.QuantumMechanics.FiniteDimensional.
    rankOneDefectMatrix_mulVec_rankOneResolventVector_eq_smul
-- [propext, Classical.choice, Quot.sound]
\end{verbatim}
The repository admission audit reports no repository-owned \texttt{axiom}, \texttt{sorry}, or
\texttt{admit} for this declaration.

\section{AI, automation, and review record}
The Lean proof and this exposition were AI-assisted and checked against the source and theorem
manifest by the maintainer.  Lean kernel checking, build output, hashing, and LaTeX compilation are
automation, not independent expert review.  No independent mathematical, scientific, novelty, or
publication review is recorded.

\section{Scientific, numerical, and certificate boundaries}
Without a separate nonzero-source premise and \texttt{QM-RESOLVENT-THM-012}, the result is an equation, not a nonzero eigenvector assertion.  It proves no spectral isolation, normalization, numerical root correctness, or infinite-volume result.  No runtime certificate is created by this dossier, and compilation of this
PDF proves nothing beyond successful typesetting.

\section{Unresolved questions}
A physical application must certify a nonzero vector, exact or enclosed secular zero, shifted invertibility, parameter provenance, and any desired spectral-isolation statement.

\section{Review and revision history}
\begin{tabularx}{\textwidth}{@{}lllY@{}}
\toprule
Date & Revision & Process & Outcome and limitation \\
\midrule
2026-10-26 & initial dossier & AI-assisted maintainer review & exact proposition, proof structure, and boundaries checked; no independent expert review \\
\bottomrule
\end{tabularx}

\end{document}
